\input pdfToolbox
\input pdfmsym
\pdfmsymsetscalefactor{10}
\input ../preamble
\def\enumindent{0cm}

\def\printmcount{\the\counter{section}.\the\counter{math counter}}
\setcounter{section}{7}

\centerline{\setfontandscale{bf}{20pt}Commutative Algebra \& Algebraic Geometry}
\medskip
\centerline{\bf Exercise \the\counter{section}}
\smallskip
\centerline{\it Ari Feiglin}

\def\rad{{\rm rad}}

\bprob

    \benum
        \item Let $X$ be an affine variety.
        Prove the following.
        \benum
            \item Let $\cll_X$ be the set of closed subsets of $X$, and $\rad_X$ the set of radical ideals of $k[X]$.
            We have maps $\I\colon\cll_X\to\rad_X$ and $Z\colon\rad_X\to\cll_X$; show these form an antiequivalence of posets.

            \item Show that these maps send irreducible closed subsets to prime ideals and vice versa.

            \item Show that these maps send points to maximal ideals and vice versa.
        \eenum

        \item Let $f\colon X\to Y$ be a morphism of affine varieties, inducing a map $f^*\colon k[Y]\to k[X]$.
        Prove the following.
        \benum
            \item Let $T\subseteq X$ be a subset, then $(f^*)^{-1}(\I_X(T))=\I_Y(f(T))$.
            \item $f^*$ is injective iff $f$ is {\emphcolor dominant} ($f(X)$ is dense in $Y$).
            \item $f^*$ is surjective iff $f$ is a closed embedding.
        \eenum
    \eenum

\eprob

\def\prime{{\rm prime}}

\benum
    \item For an affine variety $X\subseteq\bA^n$, we know that $X$ is homeomorphic to $\specm(A)$ for $A=k[X]$.
    The homeomorphism is given by mapping $P=(a_1,\dots,a_n)\in X$ to $\m_P=(x_1-a_1,\dots,x_n-a_n)$, call it $\phi$.

    Denote the zero set of $T\subseteq k[X]$ by $Z_X(T)=\set{P\in X}[f(P)=0\hbox{ for all $f\in T$}]$.
    And denote its vanishing set in $\specm(A)$ by $V_A(T)=\set{\m\in\specm(A)}[T\leq\m]$.

    Now we claim that $\phi$ maps $Z_X(T)$ to $V_A(T)$.
    Indeed, we know $f(P)=0$ iff $f\in\m_P$, so $P\in Z_X(T)$ iff $f(P)=0$ for all $f\in T$ iff $f\in\m_P$ for all $f\in T$
    iff $T\subseteq\m_P$ iff $\phi(P)\in V_A(T)$.
    That is,
    $$ \phi(Z_X(T)) = V_A(T) . $$
    This shows that $\phi$ is a closed map, and since it is bijective (as $\m_P$ exhaust the maximal ideals of $A$), its inverse is closed as well.
    Thus it forms a homeomorphism.

    We also have functions $\I_X$ and $\I_A$ from $\cll_X$ and $\cll_{\specm(A)}$ respectively to $\rad_A$ given by
    $$ \I_X(S) = \set{f\in A}[f(P)=0\hbox{ for all $P\in S$}],\qquad \I_A(Z) = \bigcap_{\m\in Z}\m . $$
    We claim that $\I_X(S)=\I_A(\phi(S))$.
    Indeed, $f\in\I_A(\phi(S))$ iff $f\in\m$ for all $\m\in\phi(S)$, i.e.\ iff $f\in\m_P$ for all $P\in S$.
    As before, this is equivalent to $f(P)=0$ for all $P\in S$, which is equivalent to $f\in\I_X(S)$.

    So we have the following diagram where each triangle commutes

    \bigskip
    \centerline{\drawdiagram{
        $\cll_X$&&$\cll_{\specm(A)}$\cr
        &$\rad_A$
    }{
        \diagarrow{from={1,1}, to={1,3}, text=$\phi$, y distance=.25cm}
        \diagarrow{from={1,1}, to={2,2}, text=$\I_X$, x distance=-.4cm, x off=-.1cm}
        \diagarrow{from={2,2}, to={1,1}, text=$Z_X$, x distance=.3cm, x off=.1cm}
        \diagarrow{from={1,3}, to={2,2}, text=$\I_A$, x distance=.4cm, x off=.1cm}
        \diagarrow{from={2,2}, to={1,3}, text=$V_A$, x distance=-.5cm, x off=-.1cm}
    }}

    Now we know that
    $$ \I_A\circ V_A(I) = \bigcap_{\m\in V_A(I)}\m = \bigcap_{I\leq\m}\m = \jac_A(I) = \sqrt I = I , $$
    where the second-to-last equality is because $A$ is Jacobson.
    Now consider $\I_A^\prime$ and $V_A^\prime$ which are the ideal and vanishing set functions for $\spec(A)$, i.e.\ we have
    $$ \I_A^\prime(S) = \bigcap_{\p\in S}\p,\qquad V_A^\prime(T) = \set{\p\in\spec(A)}[T\subseteq\p] . $$
    We know that $V_A^\prime\circ\I_A^\prime(S)=\cl S$ in $\spec(A)$.
    Clearly for $S\subseteq\specm(A)$ we have $\I_A^\prime(S)=\I_A(S)$ since they are both just the intersection of $S$.
    Furthermore $V_A(\bul)=V_A^\prime(\bul)\cap\specm(A)$, thus
    $$ V_A\circ\I_A(S) = (V_A^\prime\circ\I_A^\prime(S))\cap\specm(A) = \cl(S)\cap\specm(A) = \cl(S) , $$
    where the first closure is in $\spec(A)$, and the equality is because closure in a subspace is just the intersection of closure in the larger
    space with the subspace.
    So for $S$ closed this is the identity.

    Thus $\I_A,V_A$ are inverses.
    And since they are both isomorphic to $\I_X$ and $Z_X$ we get the desired result:
    $$ \I_X\circ Z_X = (\I_A\circ\phi)\circ(\phi^{-1}\circ V_A) = \I_A\circ V_A = \id, $$
    and
    $$ Z_X\circ\I_X = (\phi^{-1}\circ V_A)\circ(\I_A\circ\phi) = \phi^{-1}\circ\phi = \id $$
    as desired.

    Since $\I_A,V_A$ are order-reversing so too are $\I_X=\I_A\circ\phi,Z_X=\phi^{-1}\circ V_A$ (since $\phi,\phi^{-1}$ preserve order
    as bijections).
    This completes the proof.

    \item From the previous homework we know that $Z\subseteq\spec(A)$ closed is irreducible iff $\abs Z\subseteq\specm(A)$ is closed
    irreducible, and $\abs Z=Z\cap\specm(A)$.
    All closed subsets of $\specm(A)$ are of this form (by the subspace topology).

    We showed in a previous homework that the closed irreducible subsets of $\spec(A)$ are those of the form $V_A^\prime(\p)$ for $\p\leq A$
    prime.

    So if $Z\subseteq\specm(A)$ is irreducible, then it is of the form $Z=\abs Y$ for $Y\subseteq\spec(A)$ irreducible (since $Y$ is irreducible
    iff $\abs Y$ is).
    And so $Y=V_A^\prime(\p)$ for prime $\p$, and so $Z=V_A(\p)$ and since $\I_A$ is inverse to $V_A$,
    $$ \I_A(Z) = \p . $$

    Since $\phi$ is a homeomorphism it maps closed irreducible sets to closed irreducible sets and vice versa.
    So for $Z\subseteq X$ closed irreducible, $\phi(Z)\subseteq\specm(A)$ is closed irreducible and so as we just showed
    $$ \I_X(Z) = \I_A\circ\phi(Z) = \p $$
    for some prime $\p$ as desired.

    Conversely for $\p$ prime, $V_A^\prime(\p)$ is irreducible and thus $V_A(\p)$ is irreducible in the subspace topology.
    Then $Z_X(\p)=\phi^{-1}V_A(\p)$ is irreducible since $\phi$ is a homeomorphism.

    \item A point of $\specm(A)$ is a maximal ideal $\m$, and we have $V_A(\m)=V_A^\prime(\m)=\set m$ by maximality.
    And then
    $$ \I_A(\set\m) = \I_A\circ V_A(\m) = \m , $$
    so $\I_A$ sends points to maximal ideals.

    Then
    $$ \I_X(\set P) = \I_A\circ\phi(\set P) $$
    is also a maximal ideal since $\phi(\set P)$ is a point.
    And
    $$ Z_X(\m) = \phi^{-1}\circ V_A(\m) = \phi^{-1}\set\m $$
    is also a point.

    \item We view $X,Y$ as closed subsets of affine space $\bA^n,\bA^m$.
    Let $A=k[X]$ and $B=k[Y]$, then $f^*$ induces a map $\tilde f\colon\specm(A)\to\specm(B)$ which maps a maximal ideal $\m\leq A$
    to $(f^*)^{-1}(\m)\leq B$.
    Let $\phi\colon X\to\specm(A)$ and $\psi\colon Y\to\specm(B)$ be the homeomorphisms mapping $P\mapsto\m_P$.
    Then we have the following diagram

    \commsq{\specm(A)&\specm(B)\cr X&Y}
        {\tilde f}{f}{\phi^{-1}}{\psi^{-1}}

    We claim this commutes, so
    $$ \psi\circ f = \tilde f\circ\phi . $$
    Indeed, let $P\in X$ then
    $$ \tilde f\circ\phi(P) = \tilde f(\m_P) = \set{q\in B}[f^*(q)\in\m_P],\qquad
    \psi\circ f(P) = \m_{f(P)} . $$
    Since everything is well-defined here, both of these are maximal ideals.
    So to show they are equal it is sufficient to show that one is contained in the other.

    For $q\in B$ such that $f^*(q)\in\m_P$ we have that $f^*(q)=q\circ f$ is in $\m_P$, meaning $q(f(P))=0$.
    Since $\m_{f(P)}$ is the maximal ideal consisting of polynomials where $f(P)$ is a root, we see that $q\in\m_{f(P)}$.
    So we have inclusion and thus equality.

    Now,
    $$ \I_Y(f(T)) = \I_B(\psi(f(T))) = \bigcap_{\m\in\psi(f(T))}\m . $$
    On ther other hand,
    $$ \eqalign{
        (f^*)^{-1}(\I_X(T)) &= (f^*)^{-1}(\I_A\circ\phi(T))\cr
            &= (f^*)^{-1}\bigcap_{\m\in\phi(T)}\m\cr
            &= \bigcap_{\m\in\phi(T)}(f^*)^{-1}\m\cr
            &= \bigcap_{\m\in\phi(T)}\tilde f(\m)\cr
            &= \bigcap_{\m\in\phi(T)}\psi\circ f\circ\phi^{-1}(\m)\cr
            &= \bigcap_{x\in T}\psi\circ f(x),
    } $$
    and this is just the intersection of $\psi(f(T))$, so $\I_Y(f(T))$ as desired.

    \item Let $\tilde g\colon\spec(A)\to\spec(B)$ be the map induced by $f^*\colon B\to A$, which restricts to $\tilde f$ on the maximum
    spectra.
    We showed that $\tilde g$ is dominant iff $\ker f^*\subseteq\nil(B)$, and since $B$ is reduced this is if $f^*$ is injective.
    We claim that $\tilde g$ is dominant iff $\tilde f$ is.

    Now we claim that
    $$ \tilde f(\specm(A)) = \tilde g(\spec(A))\cap\specm(B) = \abs{\tilde g(\spec(A))} . $$
    The left-to-right inclusion is clear.
    Now if $\tilde g(\p)$ is maximal, then for $\p\leq\m$ maximal we have $\tilde g(\p)\leq\tilde g(\m)$ and by maximality this must be
    maximal.
    So $\tilde g(\p)=\tilde g(\m)=\tilde f(\m)$, so the two sets are equal.

    Now we claim that for $S\subseteq\spec(A)$, if $\abs S\subseteq\specm(A)$ is dense then $S$ is.
    We know that
    $$ \cl{\abs S} = V_A\circ\I_A(\abs S), $$
    and previously we showed that $V_A(I)=\specm(A)$ iff $V_A^\prime(I)=\spec(A)$.
    And $V_A^\prime(I)=\spec(A)$ iff $I\subseteq\nil(A)=0$, so iff $I=0$.
    So $\abs S$ is dense iff $\I_A(\abs S)=0$, and $\I_A(S)\subseteq\I_A(\abs S)$ so this implies $S$ is dense.

    So if $\tilde f$ is dominant, then so is $\tilde g$ and so $f^*$ is injective.
    And conversely if $f^*$ is injective then $\tilde g$ is dominant and in $\spec(A)$,
    $$ \cl{\tilde f(\specm(A))} = \cl{\tilde g(\specm(A))} \supseteq \tilde g(\cl{\specm(A)}) = \tilde g(\spec(A)) . $$
    This is dense and closed, so it must be all of $\spec(A)$.
    So $\tilde f(\specm(A))$ is dense in $\spec(A)$ and contained in $\specm(A)$, so it is dense in $\specm(A)$.
    Thus $\tilde f$ is dominant when $f^*$ is injective.

    Here we used that $\specm(A)$ is dense in $\spec(A)$.
    This is because its closure is
    $$ V_A^\prime\parens{\bigcap_{\m\in\specm(A)}\m} = V_A^\prime(\jac(A)) = \spec(A) , $$
    since the Jacobson is equal to the nilradical as $A$ is Jacobson.

    So we see that $f^*$ is injective iff $\tilde f$ is dominant.
    Since $f$ and $\tilde f$ are topologically conjugate, $\tilde f$ is dominant iff $f$ is so we conclude the proof.

    \item As before, we prove this for $\tilde f$ and appeal to the fact that $f$ and $\tilde f$ are topologically conjugate.

    If $f^*$ is surjective then $\tilde g$ is a closed embedding.
    Then its restriction $\tilde f$ is a closed embedding since it is a homeomorphism onto its image as a restriction of a homeomorphism,
    and its image is $\tilde f(\specm(A))=\tilde g(\spec(A))\cap\specm(B)$ as above, which is relatively closed.
    So if $f^*$ is surjective then $\tilde f$ and thus $f$ is a closed embedding.

    Conversely if $f$ is a closed embedding, then $f(X)\subseteq Y$ is a closed subset isomorphic to $X$, thus an isomorphic affine variety.
    Inclusion $f(X)\inj Y$ is a morphism which induces a map $k[Y]\to k[f(X)]$ which is surjective.
    This is because $k[f(X)]$ are polynomial functions $f(X)\to k$ and thus the restriction of a polynomial function $Y\to k$.
    Then the composition
    $$ k[Y] \xtos{\inc^*} k[f(X)] \xtos{f^*} k[X] $$
    is surjective since $f$ is an homeomorphism between $X$ and $f(X)$, and by functorality it is equal to $(\inc\circ f)^*=f^*$.
\eenum

\bprob

    Fix a variety $X$.
    \benum
        \item Let $X=\bP^2$ and let $P=[a:b:c]\in\bP^2$, compute the stalk $\O_{\bP^2,P}$.
        \item Give a criterion on a germ $(U,f)\in\O_{X,p}$ to be a unit.
        \item Assume $X$ is irreducible, let $p\in X$ and $U,V$ be open neighborhoods of $p$, $f\in\O_X(U)$ and $g\in\O_X(V)$ such that
        the germs $(U,f)$ and $(V,g)$ are equal in $\O_{X,p}$.
        \item Find a counterexample to the above point if $X$ is not irreducible.
    \eenum

\eprob

\benum
    \item In recitation we showed that for an affine neighborhood $p\in U$ we have
    $$ \O_{X,p} \cong k[U]_{\m_P} $$
    where $\m_P$ is the maximal ideal of functions zero at $P$.

    Suppose $P=[a:b:1]$, the other cases (when $c=0$) are solved similarly.
    $P$ has an affine neighborhood isomorphic to $\bA^2$, and $P$ is isomorphic to the point $(a,b)$.
    Then
    $$ \O_{\bP^2,P} \cong k[\bA^2]_{(x_1-a,x_2-b)} . $$
    The isomorphism $k[\bA^2]\to k[\bA^2]$ obtained by mapping $x_1\mapsto x_1-a$ and $x_2\mapsto x_2-b$ (with inverse obtained by adding $a$
    and $b$) lifts to an isomorphism $k[\bA^2]_{(x_1-a,x_2-b)}\cong k[\bA^2]_{(x_1,x_2)}$.
    So in general
    $$ \O_{\bP^2,P} \cong k[x_1,x_2]_{(x_1,x_2)} . $$

    \item We claim that $(U,f)$ is a unit iff $f(P)\neq0$.
    If $(U,f)$ is a unit there is a germ $(V,g)$ such that $(U,f)\cdot(V,g)=(U\cap V,f\cdot g)$ is the identity, so equivalent to the germ of 
    the constant $1$.
    Thus $f\cdot g(P)=1$, which means $f(P)\neq0$.

    Conversely if $f(P)\neq0$ then $P\in D_U(f)$ and so $(D_U(f),1/f)$ is a germ in $\O_{X,P}$.
    And $(U,f)\cdot(D_U(f),1/f)=(D_U(f),1)$ is a germ of $1$, so it is the identity.
    Thus $(U,f)$ is a unit.

    \item Irreduciblity is equivalent to all nonempty open subsets having nonempty intersection, equivalently being dense.
    This means an open subset of an irreducible space is irreducible (open subsets of an open subset are open).

    By definition $f|_W=g|_W$ on some subset $W\subseteq U\cap V$.
    Since $X$ is irreducible, $U\cap V$ is irreducible and thus $W$ is dense in $U\cap V$.
    By the universal property of $\bA^1$, regular maps in $\O_X(U)$ are continuous $U\to\bA^1$.
    Thus the map $f-g\colon U\cap W\to k$ is continuous, and so its zero set $Z(f-g)$ is closed in $U\cap V$ (as the proper closed sets of $\bA^1$
    are the finite sets, so $\set0$ is closed).
    Since it contains a dense set $W$, it must be all of $U\cap V$ (since it contains $\cl W=U\cap V$).
    Thus $f=g$ on $U\cap V$.

    \item Consider the disjoint union $X=\bA^1\amalg\bA^1$.
    This has an open cover of $\bA^1$s, so it is a variety.
    And by the universal properties of affine space and the coproduct
    $$ \O_X(X) \cong \hom(X,\bA^1) \cong \hom(\bA^1,\bA^1)^2 \cong k[x]^2 . $$
    So a regular map $f\in\O_X(X)$ is given by two polynomials $f_1,f_2$ which dictate $f$'s behavior on the two copies of the affine line.

    So let $p,q,r\in k[x]$ be three distinct polynomials.
    Define $f$ to be the regular function given by $(p,q)$, and $g$ be the regular function $(p,r)$.
    Take a point $P$ in the first copy of the affine line, so $(X,f)$ and $(X,g)$ both restrict to $p$ on an open neighborhood of $P$ (the first
    copy of the affine line).
    But clearly $f$ and $g$ don't agree on all of $X$.
\eenum

\def\pdv#1{\frac{}{\partial#1}}
\def\pdvof#1#2{\frac{\partial#1}{\partial#2}}
\def\ipdvof#1#2{{\partial#1}/{\partial#2}}
\bprob

    For coprime $f,g\in k[x,y]$, let $C=Z(f)$ and $D=Z(g)$ be the curves they define.
    \benum
        \item Let $f=x^2+y^2-1$ and $g=y-x^2+1$.
        Determine the intersection of $C$ and $D$ and for which points this intersection is transversal.
        \item Let $\bar f,\bar g$ be the linear parts of the Taylor expansions of $f$ and $g$, respectively, centered around $P\in C\cap D$.
        Show that $f$ and $g$ intersect transversally iff $\bar f,\bar g$ form a basis of $\m_P/\m_P^2$.
        \item Define the {\emphcolor Jacobian} of $f$ and $g$ to be the polynomial obtained as the determinant of the matrix of their partial derivatives:
        $$ J(f,g) = \det\pmatrix{\pdvof fx&\pdvof fy\cr\pdvof gx&\pdvof gy} . $$
        Show that for $P\in C\cap D$, $C$ and $D$ intersect transversally at $P$ iff $J(f,g)\neq0$.
    \eenum

\eprob

\benum
    \item We know
    $$ C\cap D = Z(f,g) , $$
    so we first simplify the ideal $(f,g)$.
    In our quotient we have $y=x^2-1$, so we can subsitute that in $f$ to get
    $$ x^2 + (x^2-1)^2 - 1 = x^4 - x^2 , $$
    thus
    $$ (f,g) = (y-x^2+1,x^2(x^2-1)) . $$
    So
    $$ C\cap D = Z(y-x^2+1)\cap Z(x^2(x^2-1)) . $$
    Zeroes of the latter term are given by $x=0$ or $x=\pm1$.
    As $y=x^2-1$ these give the intersection points
    $$ (0,-1),\qquad (1,0),\qquad (-1,0) . $$

    The stalk at the point at $P=(a,b)$ is $k[x,y]_{(x-a,y-b)}$.
    Notice that
    $$ (y+x^2-1)g = y^2 - x^4 + 2x^2 - 1 $$
    so we have
    $$ f - (y+x^2-1)g = x^2(x^2-1) . $$
    Since $x^2$ is not in $(x\pm1,y)$ when localizing by these maximal ideals it is a unit so we see that
    $$ x^2 - 1 \in (f,g)k[x,y]_{(x\pm1,y)} , $$
    and thus $x\pm1$ is in the right-hand ideal (since $x\mp1$ is a unit).
    Similarly
    $$ f + g = y(1+y) , $$
    and since $1+y$ is a unit in the localization by $(x\pm1,y)$ we see that $y$ is in the ideal as well.
    Thus in the localization $(x\pm1,y)\subseteq(f,g)$ and since the right-hand ideal is proper (as $(f,g)\subseteq(x\pm1,y)$) the ideal is maximal.
    So $f,g$ intersect transversally at $(\pm1,0)$.

    Our calculations show that $1+y$ is in the localization of $(f,g)$ by $(x,y+1)$ as well, but we claim that the intersection is not transversal
    If it were, then $C$ and $D$ would intersect transversally at all the points in their intersection so $(f,g)$ would be radical.
    But as above we know $x^2(x^2-1)\in(f,g)$ and so $x^2(x^2-1)^2\in(f,g)$, and if it were radical then $x(x^2-1)\in(f,g)$.

    But if we quotient $k[x,y]$ by $(y-x^2+1)$ we get $k[x]$ (by mapping $y$ to $x^2-1$).
    Then $(f,g)$ corresponds to $(x^2(x^2-1))$ and $x(x^2-1)$ to itself.
    And clearly $x(x^2-1)$ is not $(x^2(x^2-1))$, a contradiction.

    \item Denote $P=(P^0,P^1)$.
    Notice that every polynomial $f\in k[x,y]$ has an expansion around $P$.
    This is because the morphism $k[x,y]\mapsto k[x,y]$ given by $x\mapsto x-P^0$ and $y\mapsto y-P^1$ is an isomorphism, given by $x\mapsto x+P^0,y\mapsto y+P^1$.
    So $f\in k[x,y]$ is equal to the image of some polynomial $\sum_{n,m}a_{nm}x^ny^m$, i.e.
    $$ f = \sum_{n,m}a_{n,m}(x-P^0)^n(y-P^1)^m . $$
    For such a polynomial, we have $f(P)=a_{0,0}$ and by linearity and the product rule we have
    $$ \pdvof fx(P) = a_{1,0},\qquad \pdvof fy(P) = a_{0,1} . $$

    Let us define a map $\m_P\to\bA^2$ by $h\mapsto(\ipdvof hx,\ipdvof hy)(P)$.
    This is $\bR$-linear by linearity of the derivative.
    Furthermore it is surjective: for $(a,b)\in\bA^2$ define $h=a(x-P^0)+b(y-P^1)$ then we see that $h(P)=0$ and $\ipdvof hx=a,\ipdvof hy=b$.

    We also show that its kernel is $\m_P^2$.
    For $p,q\in\m_P$ we see that
    $$ \pdv x(pq) = \parens{\pdvof px}q + p\parens{\pdvof qx} , $$
    and evaluating at $P$ gives $0$ as $p(P)=q(P)=0$; similarly for $y$.
    So $\m_P^2$ is contained in the kernel.
    Conversely for $f\in\m_P$ in the kernel we can consider its expansion around $P$, $f=\sum_{n,m}a_{n,m}(x-P^0)^n(y-P^1)^m$.
    It being in the kernel means that $a_{0,0}=a_{1,0}=a_{0,1}=0$, and so all of the terms in $f$ are of the form $(x-P^0)^n(y-P^1)^m$ for $n+m\geq2$ (not both $\leq1$)
    which is then in $\m_P^2$.

    So this map quotients to a linear isomorphism $\m_P/\m_P^2\cong\bA^2$.
    Via this isomorphism, we see that $(1,0),(0,1)\in\bA^2$ have preimages $x-P^0$ and $y-P^1$ (from our previous computation) respectively.
    Furthermore for a polynomial $f\in k[x,y]$ with expansion around $P$ given by the coefficients $a_{n,m}$, by definition we have
    $$ \bar f = a_{0,0} + a_{1,0}(x-P^0) + a_{0,1}(y-P^1),\qquad a_{0,0} = f(P),\quad a_{1,0} = \pdvof fx(P),\quad a_{0,1} = \pdvof fy(P) . $$
    For $f\in\m_P$ this is just its projection into $\m_P/\m_P^2$.

    Now, we know that $\O_{\bA^2,P}=k[x,y]_{\m_P}$ which is a local ring with maximal ideal $\m_Pk[x,y]_{\m_P}$ and $\m_P=(x-P^0,y-P^1)$.
    By a change of coordinates, we can assume $P^0=P^1=0$ (via mapping $x\mapsto x-P^0$ and $y\mapsto y-P^1$).
    This doesn't change any of the proof, but it does tidy it up.
    So $C,D$ intersect transversally at $P=(0,0)$ iff $x,y\in(f,g)k[x,y]_{\m_P}$.
    Let $a_{n,m}$ be the coefficients of $f$, and $b_{n,m}$ the coefficients of $g$.

    In such a case, after multiplying out the denominators, we see that there are polynomials $p,q,r\in k[x,y]$ with $p\notin\m_P$ such that
    $$ xp = fq + gr . $$
    To show that $\bar f,\bar g$ form a basis of $\m_P/\m_P^2$ it is sufficient to show that $x$ and $y$ are in their span, as they form a basis for $\m_P/\m_P^2$.
    Differentiating the above identity and then evaluating at $P$ we see
    $$ p(P) = \pdvof fx(P)q(P) + \pdvof gx(P)r(P) = a_{1,0}q(P) + b_{1,0}r(P) . $$
    Let $c=p(P)$, $\alpha=q(P)$, and $\beta=r(P)$; $c\neq0$ since $p\notin\m_P$.
    Then $c=\alpha a_{1,0}+\beta b_{1,0}$ so
    $$ cx = \alpha a_{1,0}x + \beta b_{1,0}x = \alpha\bar f + \beta\bar g $$
    since $a_{0,0}=b_{0,0}=0$ as $f,g\in\m_P$.
    Since $c$ is nonzero and thus a unit, this shows that $x$ is in the span of $\bar f,\bar g$ as desired, and similarly so is $y$.
    Thus if $f,g$ intersect transversally at $P$ then $\bar f,\bar g$ form a basis.

    Conversely suppose $\bar f,\bar g$ form a basis, we want to show they intersect transversally.
    It is sufficient to show that $x,y$ are in $(f,g)$ in the stalk $\O_{\bA^2,P}=k[x,y]_{\m_P}$ as its maximal ideal is generated by $x,y$ (and $(f,g)$ is proper since it
    is contained in $(x,y)=\m_P$).
    Since $\bar f,\bar g$ form a basis there are $\alpha,\beta$ scalars such that $x=\alpha\bar f+\beta\bar g$.
    Since $f-\bar f=\tilde f\in\m_P^2$ we have that
    $$ x = \alpha(f-\tilde f) + \beta(g-\tilde g) = \alpha f + \beta g - (\alpha\tilde f + \beta\tilde g) \in (f,g) + \m_P^2 . $$
    And similarly $y\in(f,g)+\m_P^2$ so that
    $$ \m_P = (f,g) + \m_P^2 . $$
    By Nakayama's lemma (the second formulation, with $I=M=\m_P$ in the local ring $k[x,y]_{\m_P}$ and $N=(f,g)$) we see $(f,g)=\m_P$ as desired (all in the stalk).

    \item Via the isomorphism $\m_P/\m_P^2\to\bA^2$, $\bar f=a_{1,0}x+a_{0,1}y,\bar g=b_{1,0}x+b_{0,1}$ form a basis iff $(a_{1,0},a_{0,1})$ and $(b_{1,0},b_{0,1})$ do in $\bA^2$.
    This is iff the determinant of the matrix they form has nonzero determinant.
    Since $a_{1,0}=\ipdvof fx(P)$ and $a_{0,1}=\ipdvof fy(P)$ we get the desired result.
\eenum

\bprob

    Let $X$ be an irreducible variety.
    \benum
        \item Let $U$ be a nonempty affine open subset of $X$.
        Show that $\O_X(U)$ is a domain.
        Further show that $\Frac(\O_X(U))$ is independent, up to canonical isomorphism, on the choice of $U$.
        Define the {\emphcolor function field} of $X$ to be $k(X)=\Frac(\O_X(U))$.

        \item Compute $k(\bP^2)$.

        \item Define embeddings $\O_X(X)\to k(X)$ and $\O_{X,P}\to k(X)$ for all $P\in X$.

        \item Show that the image of $\O_{X,P}$ in $k(X)$ are $f/g\in\Frac(\O_X(U))$ where $U$ is an affine open neighborhood of $P$
        and $g(P)\neq0$.

        \item Prove that $\O_X(X)=\bigcap_P\O_{X,P}$ in $k(X)$.
    \eenum

\eprob

\benum
    \item Since $X$ is irreducible, its open subsets are irreducible as well.
    Thus $U$ is an irreducible affine variety, and we showed in question 1 that this means $\I(U)$ is prime.
    Then we know $\O_X(U)\cong k[U]=k[\bA^n]/\I(U)$ which is then a domain.

    We claim that $k(X)$ is equal to the filtered colimit $\colim_{V\subseteq X}\O_X(V)$ over all $V\subseteq X$ nonempty open.
    Then clearly it is independent on choice of affine neighborhood.
    For any nonempty open $U\subseteq X$ by irreducibility $U$ is dense.
    Let $\Op_*(X)$ denote the nonempty open subsets of a topological space $X$.
    For $X$ irreducible this is cofiltered, as the intersection of any two nonempty open subsets is nonempty.
    And $\Op(U)_*^\op\subseteq\Op(X)_*^\op$ is cofinal: for every $V\subseteq X$ nonempty, $U\cap V\subseteq V$ is a nonempty neighborhood
    of $U$.
    Since filtered colimits can be taken over cofinal subsets we see that for any $U\subseteq X$ open,
    $$ \colim_{V\subseteq X}\O_X(V) \cong \colim_{V\subseteq U}\O_X(V) . $$

    So now it remains to show that the right-hand colimit is isomorphic to $\Frac(\O_X(U))$.
    This is just a question about affine varieties.
    We see that elements of the colimit are elements $(V,f)$ for $f\in\O_X(V)$ where we identify $(V,f)$ and $(W,g)$ if they restrict to
    the same function on some common open subset of $V$ and $W$.
    Using the same argument as for stalks, this is iff $f|_{V\cap W}=g|_{V\cap W}$.

    For the sake of conciseness we will define $k(X)$ to be the colimit and show that it is isomorphic to $\Frac(\O_X(U))$.

    There is a canonical morphism for all open $V\subseteq X$ from $\O_X(V)$ to $k(X)$ (as it is the colimit) given by $f\mapsto(V,f)$.
    In particular this is true for $U$, and we claim that it maps nonzero elements to units of $k(X)$.
    For $f\in\O_X(U)$ nonzero, $D_U(f)$ is then nonempty and so there is an element $(D_U(f),1/f)\in k(X)$.
    This an inverse for $(U,f)$, so it is a unit.
    By localization, this induces a morphism which we see is
    $$ \phi\colon\Frac(\O_X(U)) \to k(X),\qquad f/g\mapsto (D_U(g),f/g) . $$
    This map is injective: if $\phi(f/g)=0$ then $f/g$ is identitfied with $0$ and thus $f/g|_{D_U(g)}=0$ implying $f|_{D_U(g)}=0$.
    By density again this means $f=0$ on $U$ so $f/g=0$.

    It is also surjective.
    For $V\subseteq U$ open we know that $\O_X(V)$ are rational functions on $V$, say $f/g$ for $0\notin g(V)$ for polynomials $f,g$.
    Then $f,g\in\O_X(U)$ as they are polynomial, $g\neq0$ so $f/g$ is an element of $\Frac(\O_X(U))$.
    Thus we have an isomorphism
    $$ k(X) \cong \Frac(\O_X(U)) . $$

    \item $\bP^2$ has an affine neighborhood $\bA^2$, so
    $$ k(\bP^2) = \Frac(k[\bA^2]) = k(x,y) . $$

    \item There is a canonical morphism $\O_X(X)\to\colim_{U\subseteq X}\O_X(U)=k(X)$ given by mapping $f\in\O_X(X)$ to the germ $(X,f)$.
    This is injective since if $(X,f)=0$ then $f|_U=0$ for some nonempty open $U$.
    But $U$ is then dense and as before this means $f=0$.

    We have
    $$ \O_{X,P} = \colim_{P\in U\subseteq X}\O_X(U),\qquad k(X) = \colim_{U\subseteq X}\O_X(U) , $$
    so the stalk is the colimit of the restriction $\O_X$ to a subcategory of $\Op_*(X)^\op$.
    In general if $F$ is the restriction of a functor $G$, then there is a canonical map $\colim F\to\colim G$ obtained from the canonical maps
    $Fc\to\colim G$.
    In our case this map is $(U,f)\mapsto(U,f)$ (the former germ is in $\O_{X,P}$ the latter is in $k(X)$).
    This is again injective for the same reason.

    As before we can restrict our view to neighborhoods contained in an affine neighborhood (or any nonempty neighborhood) $U$.
    So our map $\O_{X,P}\to k(X)$ can be taken relative to open subsets in $U$.
    Also our map $\O_X(X)\to k(X)$ is equal to first restricting and then taking the map $\O_X(U)\to k(X)$.
    This is because the germ of $f\in\O_X(X)$ is equal to the germ of $f|_U\in\O_X(U)$ by definition.

    \item The image of $\O_{X,P}$ consists of all germs of $f\in\O_X(V)$ for $P\in V\subseteq U$ open.
    For any $f/g\in\Frac(\O_X(U))$ such that $g(P)\neq0$ we see that $f/g\in\O_X(D_U(g))$ and $P\in D_U(g)$.
    Thus $f/g$ lies in the image of $\O_{X,P}$.

    Conversely for a map in $\O_X(V)$ for $P\in V\subseteq U$ open we know that it is of the form $f/g$ for $f,g$ polynomial and
    $g$ being nonzero on $V$, in particular $g(P)\neq0$.
    So every element in the image of $\O_{X,P}$ is of the form $f/g$ for $g(P)\neq0$ and vice versa.

    \item The image of $\O_X(X)$ in $k(X)$ are all germs of functions in $\O_X(X)$.
    Clearly this is contained in the image of every $\O_{X,P}$ as per the previous point.

    If $f/g\in\bigcap_P\O_{X,P}$ then $g(P)\neq0$ for all $P$ so $g$ is nonzero on all of $X$.
    By definition for affine $U\subseteq X$, $\O_X(U)=\O_\reg(U)$ consists of all locally rational functions, so $f/g$ is in $\O_X(U)$.
    That means there is $q_U\in\O_X(U)$ such that $(U,q_U)=f/g$.
    For any two affine $U,V$ we have that $q_U$ and $q_V$ define the same germ, and by density this means $q_U|_{U\cap V}=q_V|_{U\cap V}$.
    Since $X$ has an open cover of affine subsets, by the sheaf axiom this means there is a $q\in\O_X(X)$ such that $q|_U=q_U$ for all $U\subseteq X$
    affine.
    Then clearly $q$'s germ is equal to $q_U$'s, which is $f/g$.
\eenum

\bprob

    Fix a ring $R$.
    \benum
        \item Assume $R$ is Noetherian, show that $\nil(R)^k=0$ for some $k$.
        \item Assume $R$ is Artinian, show that
        $$ R \cong \prod_{\m\in\specm(R)}R_\m . $$
    \eenum

\eprob

\benum
    \item Since $R$ is Noetherian its ideals are finitely generated, so $\nil(R)=(f_1,\dots,f_n)$ for $f_i\in R$.
    For an index set $I=(i_1,\dots,i_k)$ let $f_I=\prod_jf_{i_j}$.
    Then inductively we see
    $$ \nil(R)^k = \gen{f_I}[\hbox{$I$ is an index set of length $k$}] . $$
    Suppose that $f_i^{n_i}=0$ for all $i$, and let $k=\sum_in_i$.
    Then for every index set $I$ of length $k$, there must be some $i$ which occurs in $I$ at least $n_i$ times (pigeonhole) and then $f_I=0$.
    Thus $\nil(R)^k=0$ as desired.

    \item Since $R$ is Artinian, its max spectrum (which is equal to its spectrum) is finite, and by the previous point $\nil(R)^k=0$.
    And so
    $$ \nil(R) = \bigcap_{\m_i\in\specm(R)}\m_i = \prod_i\m_i . $$
    Thus
    $$ 0 = \nil(R)^k = \prod_i\m_i^k . $$
    The $\m_i^k$ are coprime since if $I$ is radical and contains $\m_i^k,\m_j^k$ it contains $\m_i,\m_j$ and thus is unital.
    So if $I$ contains $\m_i^k,\m_j^k$ its radical is unital and thus $I$ is unital.

    So by CRT
    $$ R \cong R/(0) \cong R/\prod_i\m_i^k \cong \prod_iR/\m_i^k . $$
    Localizing $R/\m_i^k$ at $\m_j$ is $R_{\m_j}$ quotiented by $\m_i^k$, but this is trivial if $i\neq j$ as the localizing set and $\m_i^k$
    is nontrivial.
    Rnd if $i=j$ then it is isomorphic just to $R/\m_i^k$.
    Thus by exactness
    $$ R_{\m_i} \cong \prod_j(R/\m_j^k)_{\m_i} \cong R/\m_i^k , $$
    so
    $$ R \cong \prod_iR/\m_i^k \cong \prod_iR_{\m_i} $$
    as desired
\eenum

\bprob

    Let $P_1,\dots,P_6$ be the vertices of a nondegenerate hexagon and let $L_1,\dots,L_6$ be the sides of the hexagon continued to lines.
    Let $Q_1=L_1\cap L_4$, $Q_2=L_2\cap L_5$, $Q_3=L_3\cap L_6$ be the intersection points of opposite sides.
    Write $L_i=Z_{\bA^2}(\ell_i)$ and assume that $Q_1,Q_2,Q_3$ lie on a single line $L=Z_{\bA^2}(\ell)$.
    We show that the hexagon lies on a conic.
    \benum
        \item Write $f=\ell_1\ell_3\ell_5$ and $g=\ell_2\ell_4\ell_6$, $F=Z(f)$, and $G=Z(g)$.
        Prove that $F$ and $G$ intersect transversally.

        \item Prove that $g\in(\ell,f)k[x,y]$.

        \item Deduce that $g=q\ell+cf$ for some quadratic $q\in k[x,y]$ and scalar $c\in k$.

        \item Set $C=Z(q)$ and show that $P_1,\dots,P_6\in C$.
    \eenum

\eprob

\benum
    \item Let $P\in F\cap G$ so that $f(P)=g(P)=0$.
    Without loss of generality, $\ell_1(P)=\ell_2(P)=0$.
    Since the hexagon is nondegenerate $\ell_i(P)\neq0$ for all the other $i$, so that they are units in $\O_{\bA^2,P}$.
    Thus in the stalk
    $$ (f,g) = (\ell_1,\ell_2) . $$
    Now $\ell_1,\ell_2$ are coprime (they share a root) linear polynomials in $x,y$ which are not scalar multiples of one another so they generate a
    maximal ideal in $k[x,y]$.
    Say $\ell_1=ax+by+c$ and $\ell_2=cx+dy+e$, then we can divide $\ell_1$ by $a$ (if it is zero then $b$ is not, so we can swap variables)
    and then subtract $\ell_2-c/a\ell_1$ to get a linear polynomial in $y$ and then subtract that from $\ell_1$ to get a linear polynomial in $x$.
    So $(\ell_1,\ell_2)\supseteq(x-a',y-b')$, which is maximal and thus they are equal.

    And so $(\ell_1,\ell_2)$'s image in the stalk is maximal, so $f,g$ intersect transversally.

    \item We know from AF+BG that $g\in(\ell,f)k[x,y]$ iff $g\in(\ell,f)\O_{\bA^2,P}$ for all $P$ in the intersection of $\ell,f$.
    The intersection of $f$ and $\ell$ are the points $Q_1,Q_2,Q_3$ and so $g(P)=0$ as well, so $g$ is not a unit in the stalk.
    And if $\ell(P)=f(P)=0$ then one $\ell_i$ must be zero but not the rest ($\ell_1,\ell_3,\ell_5$ don't intersect on $\ell$).
    Then we see $(\ell,f)=(\ell,\ell_i)$ in the stalk and this is maximal as they are both linear as before.
    Since the stalk is local, this is a maximal ideal, and $g$ is not a unit, this means $g\in(\ell,f)$ in the stalk.
    Thus $g\in(\ell,f)k[x,y]$.

    \item By the previous point $g=q\ell+cf$ for polynomials $q,c$.
    Since $g,f$ have order $3$ and $\ell$ has order $1$, so by the same reason in recitation we see that $q$ has order $2$ and $c$ is a scalar.

    \item For $P=P_1,\dots,P_6$ we have both $f(P)=g(P)=0$.
    Thus $q(P)\ell(P)=0$, but $P$ is a vertex of the hexagon and so it is not on $\ell$, so $q(P)=0$.
    Thus $P\in Z(q)$ as desired.
\eenum

\bprob

    \benum
        \item Show that the space of degree $m$ homogeneous polynomials in $n$ variables can be identified with $\bA^{k_m}$ for some $k_m$.
        \item Show that the general linear group $\gl_n$ acts on $\bA^{k_m}$ by a linear change of variables.
        Show that this is an action of varieties, so the action of every $g\in\gl_n$ is a morphism of varieties.
        \item We identify $(a_0,\dots,a_4)$ with the following degree $4$ homogeneous polynomial in two variables $\sum_i\binom40a_ix^{4-i}y^i$.
        For $\bar a\in\bA^5$ define its {\emphcolor Hankel matrix} to be
        $$ \pmatrix{a_0&a_1&a_2\cr a_1&a_2&a_3\cr a_2&a_3&a_4} , $$
        and define its {\emphcolor catalecticant}, $C(a)$, to be the determinant of the Hankel matrix.

        For $g\in\gl_2$ and $a\in\bA^5$ show that $C(g\cdot a)=\det(g)^6C(a)$.
        \item Show that $C\in k[\bA^5]$, and deduce that the Zariski closure $\cl{\gl_2(\bC)f}$ is contained in $Z(C)$.
        \item Deduce that $g$ is not in $\cl{\gl_2(\bC)f}$.
    \eenum

\eprob

\benum
    \item Consider an elementary homogeneous polynomial of degree $m$, so of the form $x_1^{k_1}\cdots x_n^{k_n}$ where $k_1+\cdots+k_n=m$.
    The number of choices of $k_1,\dots,k_n$ can be viewed as the number of ways to place $n-1$ bars among $m$ stripes (between the $i-1$th and
    $i$th bar we count the number of stars and set $k_i$ equal to that).
    Thus there are $k^n_m=\binom{m+n-1}{n-1}$ such elementary polynomials; they form a basis for the homogeneous polynomials of degree $m$.

    Let us enumerate the tuples $I$ where $\sum I=m$ by $I_1,\dots,I_{k^n_m}$.
    Define the map
    $$ \phi\colon\bA^{k^n_m} \to k[x_1,\dots,x_n]_{\homo,m},\qquad (a_1,\dots,a_{k^n_m}) \mapsto \sum_{i=1}^{k^n_m}a_i\binom m{I_i}x^{I_i} , $$
    where $\binom m{I_i}$ is the multinomial coefficient and $x^I=x_1^{i_1}\cdots x_n^{i_n}$.
    This is clearly linear and an isomorphism as the dimensions are equal and $\binom m{I_i}x^{I_i}$ form a basis as explained.

    \item We have $g\in\gl_n$ act on $k[x_1,\dots,x_n]$ by mapping the variable $x_i$ to
    $$ x_i \mapsto \sum_j g_{ij}x_j . $$
    This is the map on $k[x_1,\dots,x_n]_{\homo,1}$ represented by $g$ under the basis $x_1,\dots,x_n$.
    As this is a map on homogeneous polynomials of degree $1$ it extends to a map on those of degree $m$.

    This is clearly a group action, since $gh$'s action is represented by $gh$ (on homogeneous polynomials of degree $1$), and thus equal to
    the action of $h$ then $g$ (on homogeneous of degree $1$, thus on all homogeneous polynomials).

    For $x^I$ we see that $g$ maps it to a function polynomial in $x_1,\dots,x_n$.
    Thus through the linear isomorphism $\phi$, $g$ acts polynomially and thus is an isomorphism of varieties.

    \item The determinant of a matrix is polynomial in its inputs so $C\in k[\bA^5]$.
    And for $g\in\gl_2(\bC)$ we have $C(g\cdot f)=\det g^6C(f)$.
    The Hankel matrix of $f$ is ${\rm diag}(1,0,1)$ which has determinant zero, so $C(f)=0$.
    Thus $C(g\cdot f)=0$ and so $\gl_2(\bC)f\subseteq Z(C)$ and since $Z(C)$ is closed, $\cl{\gl_2(\bC)f}\subseteq Z(C)$ as desired.

    \item The Hankel matrix of $g$ is the matrix of zeroes except whose skew-diagonal consists of $1$s, which has determinant $-1$.
    So $C(g)=-1$ and so $g\notin Z(C)$ and thus is not in $\cl{\gl_2(\bC)f}$.
\eenum

\bye

